\chapter{1995年全国卷（理）}
\section{单选题}

\begin{problem}
已知 \(I\) 为全集，集合 \(M\)，\(N \subseteq I\)，若 \(M \cap N = N\)，则（\quad）

\choices
    {\(\overline{M} \supseteq \overline{N}\)}
    {\(M \subseteq \overline{N}\)}
    {\(\overline{M} \subseteq \overline{N}\)}
    {\(M \supseteq \overline{N}\)}
\end{problem}

\begin{answer}
C
\end{answer}

\begin{solution}
由 \(M\cap N=N\) 可知 \(N\subseteq M\). 补集运算会反转包含关系，因此 \(\overline{M}\subseteq\overline{N}\)，故选 C.
\end{solution}

\begin{problem}
函数 \(y = -\frac{1}{x + 1} \) 的图象是（\quad）

\choices
    {\choicebitmap[width=0.15\paperwidth]{img/1995/national_science/q2_optA_nolabel.png}}
    {\choicebitmap[width=0.15\paperwidth]{img/1995/national_science/q2_optB_nolabel.png}}
    {\choicebitmap[width=0.15\paperwidth]{img/1995/national_science/q2_optC_nolabel.png}}
    {\choicebitmap[width=0.15\paperwidth]{img/1995/national_science/q2_optD_nolabel.png}}
\end{problem}

\begin{answer}
B
\end{answer}

\begin{solution}
函数的定义域为 \(x\ne-1\)，且 \(y\ne0\). 当 \(x>-1\) 时，\(x+1>0\)，从而 \(y<0\)；当 \(x<-1\) 时，\(x+1<0\)，从而 \(y>0\). 因此图象的两支分别位于直线 \(x=-1\) 的右下方和左上方，渐近线为 \(x=-1\) 与 \(y=0\)，所以应选 B.
\end{solution}

\begin{problem}
函数 \(y = 4\sin \left(3x + \frac{\pi}{4}\right) + 3\cos \left(3x + \frac{\pi}{4}\right) \) 的最小正周期是（\quad）

\choices
    {\(6\pi\)}
    {\(2\pi\)}
    {\(\frac{2\pi}{3} \)}
    {\(\frac{\pi}{3} \)}
\end{problem}

\begin{answer}
C
\end{answer}

\begin{solution}
令 \(u=3x+\dfrac{\pi}{4}\)，则
\[
4\sin u+3\cos u=5\sin\left(u+\varphi\right)
\]
其中 \(\cos\varphi=\dfrac45\)，\(\sin\varphi=\dfrac35\). 因此函数关于 \(u\) 的最小正周期为 \(2\pi\)，而 \(u\) 随 \(x\) 的变化率为 \(3\)，故最小正周期 \(T\) 满足 \(3T=2\pi\)，即 \(T=\dfrac{2\pi}{3}\). 所以应选 C.
\end{solution}

\begin{problem}
正方体的全面积是 \(a^2\)，它的顶点都在球面上，这个球的表面积是（\quad）

\choices
    {\(\frac{\pi a^2}{3}\)}
    {\(\frac{\pi a^{2}}{2} \)}
    {\(2\pi a^{2}\)}
    {\(3\pi a^{2}\)}
\end{problem}

\begin{answer}
B
\end{answer}

\begin{solution}
设正方体棱长为 \(s\). 由全面积为 \(a^2\)，得 \(6s^2=a^2\). 球过正方体的所有顶点，所以球的直径等于正方体的体对角线 \(s\sqrt3\)，从而球半径 \(R=\dfrac{s\sqrt3}{2}\). 于是球的表面积为
\[
4\pi R^2=4\pi\left(\frac{s\sqrt3}{2}\right)^2=3\pi s^2=\frac{\pi a^2}{2}
\]
故选 B.
\end{solution}

\begin{problem}
如图，若图中的直线 \(l_{1}\)，\(l_{2}\)，\(l_{3}\) 的斜率分别为 \(k_{1}\)，\(k_{2}\)，\(k_{3}\)，则（\quad）

\bitmapfigure[max width=\linewidth]{img/1995/national_science/q5_fig1.png}

\choices
    {\(k_{1} < k_{2} < k_{3}\)}
    {\(k_{3} < k_{1} < k_{2}\)}
    {\(k_{3} < k_{2} < k_{1}\)}
    {\(k_{1} < k_{3} < k_{2}\)}
\end{problem}

\begin{answer}
D
\end{answer}

\begin{solution}
直线的斜率等于倾斜角的正切. 图中 \(l_1\) 向右下方倾斜，所以 \(k_1<0\)；\(l_3\) 向右上方倾斜且较平缓，\(l_2\) 向右上方倾斜且较陡，所以 \(0<k_3<k_2\). 因此 \(k_1<k_3<k_2\)，故选 D.
\end{solution}

\begin{problem}
在 \((1 - x^{3})(1 + x)^{10}\) 的展开式中，\(x^{5}\) 的系数是（\quad）

\choices
    {\(-297\)}
    {\(-252\)}
    {297}
    {207}
\end{problem}

\begin{answer}
D
\end{answer}

\begin{solution}
由二项式定理，\((1+x)^{10}\) 中 \(x^5\) 的系数为 \(\mathrm C_{10}^{5}=252\)，而 \(x^3(1+x)^{10}\) 中 \(x^5\) 的系数等于 \((1+x)^{10}\) 中 \(x^2\) 的系数，即 \(\mathrm C_{10}^{2}=45\). 所以原展开式中 \(x^5\) 的系数为
\[
\mathrm C_{10}^{5}-\mathrm C_{10}^{2}=252-45=207
\]
故选 D.
\end{solution}

\begin{problem}
使 \(\arcsin x > \arccos x\) 成立的 \(x\) 的取值范围是（\quad）

\choices
    {\(\left(0, \frac{\sqrt{2}}{2}\right]\)}
    {\(\left(\frac{\sqrt{2}}{2}, 1\right]\)}
    {\(\left[-1, \frac{\sqrt{2}}{2}\right)\)}
    {\([-1, 0)\)}
\end{problem}

\begin{answer}
B
\end{answer}

\begin{solution}
反三角函数的定义域要求 \(x\in[-1,1]\)，且
\[
\arccos x=\frac{\pi}{2}-\arcsin x
\]
因此
\[
\arcsin x>\arccos x\Longleftrightarrow \arcsin x>\frac{\pi}{4}
\Longleftrightarrow x>\sin\frac{\pi}{4}=\frac{\sqrt2}{2}
\]
结合定义域，得 \(x\in\left(\dfrac{\sqrt2}{2},1\right]\)，故选 B.
\end{solution}

\begin{problem}
双曲线 \(3x^{2} - y^{2} = 3\) 的渐近线方程是（\quad）

\choices
    {\(y = \pm 3x \)}
    {\(y = \pm \frac{1}{3} x \)}
    {\(y = \pm \sqrt{3} x \)}
    {\(y = \pm \frac{\sqrt{3}}{3} x \)}
\end{problem}

\begin{answer}
C
\end{answer}

\begin{solution}
将方程化为标准形式：
\[
\frac{x^2}{1}-\frac{y^2}{3}=1
\]
双曲线 \(\dfrac{x^2}{a^2}-\dfrac{y^2}{b^2}=1\) 的渐近线为 \(y=\pm\dfrac ba x\). 这里 \(a=1\)，\(b=\sqrt3\)，所以渐近线为 \(y=\pm\sqrt3 x\)，故选 C.
\end{solution}

\begin{problem}
已知 \(\theta\) 是第三象限角，且 \(\sin^4\theta + \cos^4\theta = \frac{5}{9}\)，那么 \(\sin 2\theta\) 等于（\quad）

\choices
    {\(\frac{2\sqrt{2}}{3}\)}
    {\(-\frac{2\sqrt{2}}{3}\)}
    {\(\frac{2}{3}\)}
    {\(-\frac{2}{3}\)}
\end{problem}

\begin{answer}
A
\end{answer}

\begin{solution}
由
\[
\sin^4\theta+\cos^4\theta
=\left(\sin^2\theta+\cos^2\theta\right)^2-2\sin^2\theta\cos^2\theta
=1-\frac12\sin^2 2\theta
\]
得 \(1-\dfrac12\sin^2 2\theta=\dfrac59\)，即 \(\sin^2 2\theta=\dfrac89\). 又因为 \(\theta\) 是第三象限角，\(\sin\theta<0\)、\(\cos\theta<0\)，所以 \(\sin2\theta=2\sin\theta\cos\theta>0\). 因此 \(\sin2\theta=\dfrac{2\sqrt2}{3}\)，故选 A.
\end{solution}

\begin{problem}
已知直线 \(l \perp\) 平面 \(\alpha\)，直线 \(m \subset\) 平面 \(\beta\)，有下面四个命题：

\begin{circlelist}
\item \(\alpha \parallel \beta \Rightarrow l \perp m\)；
\item \(\alpha \perp \beta \Rightarrow l \parallel m\)；
\item \(l \parallel m \Rightarrow \alpha \perp \beta\)；
\item \(l \perp m \Rightarrow \alpha \parallel \beta\).
\end{circlelist}

其中正确的两个命题是（\quad）

\choices
    {\circled{1}与\circled{2}}
    {\circled{3}与\circled{4}}
    {\circled{2}与\circled{4}}
    {\circled{1}与\circled{3}}
\end{problem}

\begin{answer}
D
\end{answer}

\begin{solution}
因为 \(l\perp\alpha\)，所以 \(l\) 垂直于平面 \(\alpha\) 内的任意直线.

当 \(\alpha\parallel\beta\) 时，\(l\) 也垂直于平面 \(\beta\)，因而 \(l\perp m\)，命题\(\circled{1}\)正确. 平面 \(\alpha\perp\beta\) 时，\(\beta\) 内只有垂直于交线的方向才平行于 \(l\)，而题设没有限定 \(m\) 的方向；取 \(\beta\) 内不平行于 \(l\) 的直线 \(m\)，即可知不能推出 \(l\parallel m\)，所以命题\(\circled{2}\)错误.

若 \(l\parallel m\)，则平面 \(\beta\) 内含有一条平行于 \(l\) 的直线，也就是含有一条垂直于平面 \(\alpha\) 的直线，故 \(\alpha\perp\beta\)，命题\(\circled{3}\)正确. 仅由 \(l\perp m\) 不能推出 \(\alpha\parallel\beta\)，所以命题\(\circled{4}\)错误. 正确的两个命题是\(\circled{1}\)与\(\circled{3}\)，故选 D.
\end{solution}

\begin{problem}
已知 \(y = \log_a(2 - ax) \) 在 \([0,1]\) 上是 \(x \) 的减函数，则 \(a \) 的取值范围是（\quad）

\choices
    {\((0,1)\)}
    {\((1,2)\)}
    {\((0,2)\)}
    {\([2, +\infty)\)}
\end{problem}

\begin{answer}
B
\end{answer}

\begin{solution}
要使函数在 \([0,1]\) 上有定义，必须有 \(2-ax>0\) 对所有 \(x\in[0,1]\) 成立. 由于 \(a>0\)，只需 \(2-a>0\)，即 \(0<a<2\). 当 \(0<a<1\) 时，\(2-ax\) 随 \(x\) 减小，而底数 \(a\) 小于 \(1\) 的对数函数是减函数，复合后为增函数；当 \(a>1\) 时，复合后为减函数. 又 \(a=1\) 不是对数的合法底数，所以 \(1<a<2\)，故选 B.
\end{solution}

\begin{problem}
等差数列 \(\{a_{n}\}\)，\(\{b_{n}\}\) 的前 \(n\) 项和分别为 \(S_{n}\) 与 \(T_{n}\)，若 \(\frac{S_{n}}{T_{n}} = \frac{2n}{3n+1}\)，则 \(\displaystyle\lim_{n \to \infty} \frac{a_{n}}{b_{n}}\) 等于（\quad）

\choices
    {1}
    {\(\frac{\sqrt{6}}{3} \)}
    {\(\frac{2}{3} \)}
    {\(\frac{4}{9} \)}
\end{problem}

\begin{answer}
C
\end{answer}

\begin{solution}
设两等差数列的首项和公差分别为 \(a_1\)，\(d_a\) 与 \(b_1\)，\(d_b\). 由前 \(n\) 项和公式，
\[
\frac{2a_1+(n-1)d_a}{2b_1+(n-1)d_b}=\frac{2n}{3n+1}
\]
由于题设比值对所有正整数 \(n\) 成立，交叉相乘后得到恒等式

\[
(3n+1)\bigl(2a_1+(n-1)d_a\bigr)=2n\bigl(2b_1+(n-1)d_b\bigr)
\]

展开并比较关于 \(n\) 的各次项系数，得到
\(3d_a=2d_b\)，\(2a_1-d_a=0\)，\(d_a=2\left(2b_1-d_b\right)\).
若 \(d_a=0\)，则由 \(2a_1-d_a=0\) 得 \(a_1=0\)，进而 \(S_n=0\)，与题设比值为正数矛盾，所以 \(d_a\ne0\). 因此 \(d_b=\dfrac32d_a\)，且
\(a_n=a_1+(n-1)d_a\sim nd_a\)，\(b_n=b_1+(n-1)d_b\sim nd_b\)
从而
\[
\lim_{n\to\infty}\frac{a_n}{b_n}=\frac{d_a}{d_b}=\frac23
\]
故选 C.
\end{solution}

\begin{problem}
用 1， 2， 3， 4， 5 这五个数字组成没有重复数字的三位数，其中偶数共有（\quad）

\choices
    {\(24\text{ 个}\)}
    {\(30\text{ 个}\)}
    {\(40\text{ 个}\)}
    {\(60\text{ 个}\)}
\end{problem}

\begin{answer}
A
\end{answer}

\begin{solution}
三位数的个位必须是偶数字 \(2\) 或 \(4\)，有 \(2\) 种选法. 确定个位后，从剩下的 \(4\) 个数字中选百位，再从剩下的 \(3\) 个数字中选十位，所以偶数共有
\[
2\cdot4\cdot3=24
\]
个，故选 A.
\end{solution}

\begin{problem}
在极坐标系中，椭圆的二焦点分别在极点和点 \((2c,0)\)，离心率为 \(e\)，则它的极坐标方程是（\quad）

\choices
    {\(\rho = \frac{c(1 - e)}{1 - e\cos\theta}\)}
    {\(\rho = \frac{c(1 - e^2)}{1 - e\cos\theta}\)}
    {\(\rho = \frac{c(1 - e)}{1 - e\cos\theta}\)}
    {\(\rho = \frac{c(1 - e^2)}{e (1 - e\cos\theta)}\)}
\end{problem}

\begin{answer}
D
\end{answer}

\begin{solution}
设椭圆的半长轴为 \(a\)，焦半距为 \(c\)，则 \(0<e<1\)，且 \(e=\dfrac ca\). 设椭圆上一点的极坐标为 \((\rho,\theta)\)，两焦点为 \(O\) 和 \(F(2c,0)\). 由椭圆的定义，
\[
\sqrt{(\rho\cos\theta-2c)^2+(\rho\sin\theta)^2}+\rho=2a
\]
移项后两边均非负，平方得
\[
\rho^2-4c\rho\cos\theta+4c^2=4a^2-4a\rho+\rho^2
\]
即
\[
\rho\left(a-c\cos\theta\right)=a^2-c^2
\]
代入 \(c=ae\)，得到
\[
\rho=\frac{a(1-e^2)}{1-e\cos\theta}
=\frac{c(1-e^2)}{e(1-e\cos\theta)}
\]
故选 D.
\end{solution}

\begin{problem}
如图，\(A_{1}B_{1}C_{1}-ABC\) 是直三棱柱，\(\angle BCA=90^{\circ}\)，点 \(D_{1}\)，\(F_{1}\) 分别是 \(A_{1}B_{1}\)，\(A_{1}C_{1}\) 的中点. 若 \(BC=CA=CC_{1}\)，则 \(BD_{1}\) 与 \(AF_{1}\) 所成的角的余弦值是（\quad）

\bitmapfigure[max width=0.42\linewidth]{img/1995/national_science/q15_fig1.png}

\choices
    {\(\frac{\sqrt{30}}{10}\)}
    {\(\frac{1}{2}\)}
    {\(\frac{\sqrt{30}}{15}\)}
    {\(\frac{\sqrt{15}}{10}\)}
\end{problem}

\begin{answer}
A
\end{answer}

\begin{solution}
取 \(C\) 为原点，令 \(CA\)、\(CB\)、\(CC_1\) 分别沿三个两两垂直的坐标轴，设它们的长度均为 \(s\). 则
\(A=(s,0,0)\)，\(B=(0,s,0)\)，\(A_1=(s,0,s)\)，\(C_1=(0,0,s)\).
因此
\(D_1=\left(\frac s2,\frac s2,s\right)\)，\(F_1=\left(\frac s2,0,s\right)\).
取方向向量
\(\overrightarrow{BD_1}=\left(\frac s2,-\frac s2,s\right)\)，\(\overrightarrow{AF_1}=\left(-\frac s2,0,s\right)\).
于是
\(\overrightarrow{BD_1}\cdot\overrightarrow{AF_1}=\frac{3s^2}{4}\)，\(|BD_1|=\frac{s\sqrt6}{2}\)，\(|AF_1|=\frac{s\sqrt5}{2}\).
两直线所成的角取锐角，所以其余弦为
\[
\frac{\overrightarrow{BD_1}\cdot\overrightarrow{AF_1}}{|BD_1|\,|AF_1|}
=\frac{3}{\sqrt{30}}=\frac{\sqrt{30}}{10}
\]
故选 A.
\end{solution}

\section{填空题}

\begin{problem}
不等式 \(\left(\frac{1}{3}\right)^{x^{2} - 8} > 3^{-2x}\) 的解集是\fillinblank{}.
\end{problem}

\begin{answer}
\(\{x\mid -2 < x < 4\}\)
\end{answer}

\begin{solution}
将左边化为以 \(3\) 为底的幂，得
\[
\left(\frac13\right)^{x^2-8}=3^{8-x^2}
\]
因为底数 \(3>1\)，原不等式等价于 \(8-x^2>-2x\)，即
\[
x^2-2x-8<0\Longleftrightarrow (x-4)(x+2)<0
\]
所以解集为 \(\{x\mid -2<x<4\}\).
\end{solution}

\begin{problem}
已知圆台上、下底面圆周都在球面上，且下底面过球心，母线与底面所成的角为 \(\frac{\pi}{3}\)，则圆台的体积与球体积之比为\fillinblank{}.
\end{problem}

\begin{answer}
\(\frac{7\sqrt{3}}{32}\)
\end{answer}

\begin{solution}
设球的半径为 \(R\)，圆台上底面半径为 \(r\)，高为 \(h\). 下底面过球心，所以其半径就是 \(R\). 在过球心和圆台轴线的截面内，母线与底面所成的角为 \(\frac{\pi}{3}\)，故
\(\tan\frac{\pi}{3}=\frac{h}{R-r}\)，\(h=\sqrt3(R-r)\).
上底面圆周也在球面上，所以
\[
r^2+h^2=R^2
\]
令 \(u=R-r\)，代入 \(h=\sqrt3u\)，得
\[
(R-u)^2+3u^2=R^2\Longleftrightarrow u(4u-2R)=0
\]
因为圆台有正的高，\(u=\dfrac R2\)，从而 \(r=\dfrac R2\)，\(h=\dfrac{\sqrt3 R}{2}\). 因此
\(V_{\text{圆台}}=\frac{\pi h}{3}\left(R^2+Rr+r^2\right) =\frac{7\sqrt3\pi R^3}{24}\)，\(V_{\text{球}}=\frac{4\pi R^3}{3}\).
所以
\[
\frac{V_{\text{圆台}}}{V_{\text{球}}}
=\frac{7\sqrt3}{32}
\]
\end{solution}

\begin{problem}
函数 \(y = \sin \left(x - \frac{\pi}{6} \right) \cos x \) 的最小值是\fillinblank{}.
\end{problem}

\begin{answer}
\(-\frac{3}{4}\)
\end{answer}

\begin{solution}
利用积化和差公式，
\[
\sin\left(x-\frac{\pi}{6}\right)\cos x
=\frac12\left[\sin\left(2x-\frac{\pi}{6}\right)+\sin\left(-\frac{\pi}{6}\right)\right]
=\frac12\sin\left(2x-\frac{\pi}{6}\right)-\frac14
\]
因为正弦函数的最小值为 \(-1\)，所以原函数的最小值为 \(-\dfrac12-\dfrac14=-\dfrac34\).
\end{solution}

\begin{problem}
直线 \(l\) 过抛物线 \(y^{2} = a(x + 1)\)（\(a > 0\)）的焦点，并且与 \(x\) 轴垂直，若 \(l\) 被抛物线截得的线段长为 \(4\)，则 \(a = \)\fillinblank{}.
\end{problem}

\begin{answer}
4
\end{answer}

\begin{solution}
将抛物线写成标准形式 \(y^2=4p(x+1)\)，得 \(p=\dfrac a4\)，所以焦点为 \(\left(-1+\dfrac a4,0\right)\). 过焦点且垂直于 \(x\) 轴的直线为 \(x=-1+\dfrac a4\). 代入抛物线方程，得
\(y^2=a\cdot\frac a4=\frac{a^2}{4}\)，\(y=\pm\frac a2\).
由于 \(a>0\)，截得的线段长为 \(\dfrac a2-\left(-\dfrac a2\right)=a\). 题设给出该长度为 \(4\)，所以 \(a=4\).
\end{solution}

\begin{problem}
四个不同的小球放入编号为 1， 2， 3， 4 的四个盒中，则恰有一个空盒的放法共有\fillinblank{}种（用数字作答）.
\end{problem}

\begin{answer}
144
\end{answer}

\begin{solution}
恰有一个空盒，等价于恰有三个盒子各至少放入一个小球. 先选空盒，有 \(\mathrm C_4^1=4\) 种选法. 对选定的三个盒子，四个不同小球放入其中且三个盒子都不空，利用容斥原理，放法数为
\[
3^4-\mathrm C_3^1 2^4+\mathrm C_3^2 1^4
=81-3\cdot16+3=36
\]
因此总放法数为
\[
\mathrm C_4^1\cdot36=4\cdot36=144
\]
\end{solution}

\section{解答题}

\begin{problem}
在复平面上，一个正方形的四个顶点按照逆时针方向依次为 \(Z_{1}\)，\(Z_{2}\)，\(Z_{3}\)，\(O\)（其中 \(O\) 是原点），已知 \(Z_{2}\) 对应复数 \(z_{2} = 1 + \sqrt{3} \i\). 求 \(Z_{1}\) 和 \(Z_{3}\) 对应的复数.
\end{problem}

\begin{solution}
正方形的对角线 \(Z_2O\) 与 \(Z_1Z_3\) 互相平分且垂直. 把从 \(O\) 指向 \(Z_2\) 的向量旋转 \(-\dfrac{\pi}{2}\) 后再缩短为原来的一半，可得到从正方形中心指向 \(Z_1\) 的向量；旋转 \(\dfrac{\pi}{2}\) 后再缩短为原来的一半，可得到从中心指向 \(Z_3\) 的向量. 因此
\(z_1=\frac{1-\i}{2}z_2\)，\(z_3=\frac{1+\i}{2}z_2\).
代入 \(z_2=1+\sqrt3\i\)，得
\[
z_1=\frac{1-\i}{2}(1+\sqrt3\i)
=\frac{\sqrt3+1}{2}+\frac{\sqrt3-1}{2}\i
\]
\[
z_3=\frac{1+\i}{2}(1+\sqrt3\i)
=\frac{1-\sqrt3}{2}+\frac{1+\sqrt3}{2}\i
\]
\end{solution}

\begin{problem}
求 \(\sin^2 20^\circ + \cos^2 50^\circ + \sin 20^\circ \cos 50^\circ\) 的值.
\end{problem}

\begin{solution}
设
\(u=\sin20^\circ+\cos50^\circ\)，\(v=\sin20^\circ-\cos50^\circ\).
由于 \(\cos50^\circ=\sin40^\circ\)，由和差化积公式得
\(u=\sin20^\circ+\sin40^\circ=\cos10^\circ\)，\(v=\sin20^\circ-\sin40^\circ=-\sqrt3\sin10^\circ\).
又 \(\sin20^\circ=\dfrac{u+v}{2}\)，\(\cos50^\circ=\dfrac{u-v}{2}\)，所以原式为
\[
\left(\frac{u+v}{2}\right)^2+\left(\frac{u-v}{2}\right)^2
+\frac{u+v}{2}\cdot\frac{u-v}{2}
=\frac{3u^2+v^2}{4}
\]
代入 \(u=\cos10^\circ\)、\(v=-\sqrt3\sin10^\circ\)，得
\[
\frac{3u^2+v^2}{4}
=\frac{3\cos^2 10^\circ+3\sin^2 10^\circ}{4}
=\frac34
\]
\end{solution}

\begin{problem}
如图，圆柱的轴截面 \(ABCD\) 是正方形，点 \(E\) 在底面的圆周上，\(AF \perp DE\)，\(F\) 是垂足.

\begin{enumerate}
\item 求证：\(AF \perp DB\).

\item 如果圆柱与三棱锥 \(D - ABE\) 的体积的比等于 \(3\pi\)，求直线 \(DE\) 与平面 \(ABCD\) 所成的角.
\end{enumerate}

\FigureLayoutDeclare{img/1995/national_science/q23_fig1.png}{5.8cm}{99bp 56bp 22bp 59bp}
\bitmapfigure[max width=0.34\linewidth]{img/1995/national_science/q23_fig1.png}
\end{problem}

\begin{solution}
\begin{enumerate}
\item 由于 \(DA\) 是圆柱的母线，\(DA\perp\) 底面 \(ABE\)，所以 \(DA\perp EB\). 又 \(AB\) 是底面圆的直径，且 \(E\) 在圆周上，所以 \(\angle AEB=90^\circ\)，即 \(AE\perp EB\). 直线 \(DA\) 与 \(AE\) 相交于 \(A\) 且都在平面 \(DAE\) 内，故 \(EB\perp\) 平面 \(DAE\)，从而 \(EB\perp AF\). 由已知 \(AF\perp DE\)，且 \(EB\) 与 \(DE\) 在 \(E\) 相交，所以 \(AF\perp\) 平面 \(DEB\)，进而 \(AF\perp DB\).

\item 过 \(E\) 作 \(EH\perp AB\)，垂足为 \(H\)，连接 \(DH\). 平面 \(ABCD\) 与平面 \(ABE\) 垂直，交线为 \(AB\)，而 \(EH\subset\) 平面 \(ABE\) 且 \(EH\perp AB\)，所以 \(EH\perp\) 平面 \(ABCD\). 因此 \(DH\) 是 \(DE\) 在平面 \(ABCD\) 上的射影，\(\angle EDH\) 就是直线 \(DE\) 与平面 \(ABCD\) 所成的角.

设圆柱底面半径为 \(R\). 由于轴截面是正方形，\(AB=AD=2R\)，所以圆柱体积为 \(2\pi R^3\). 三角形 \(ABE\) 以 \(AB\) 为直径，设 \(EH=h\)，则 \(S_{\triangle ABE}=R h\)，且三棱锥 \(D-ABE\) 的高为 \(AD=2R\)，故
\[
V_{D-ABE}=\frac13\cdot2R\cdot Rh=\frac23R^2h
\]
由题设体积比为 \(3\pi\)，得
\[
\frac{2\pi R^3}{\frac23R^2h}=3\pi\Longrightarrow h=R
\]
在圆 \(ABE\) 中，\(AB\) 是直径，且从圆周点 \(E\) 到直径 \(AB\) 的垂线段 \(EH\) 等于半径，故 \(H\) 是圆心，\(AH=R\). 在直角三角形 \(DAH\) 中，\(AD=2R\)，所以 \(DH=\sqrt5R\). 在直角三角形 \(DEH\) 中，
\[
\tan\angle EDH=\frac{EH}{DH}=\frac1{\sqrt5}
\]
故所求角为 \(\arctan\dfrac1{\sqrt5}\)，即 \(\operatorname{arccot}\sqrt5\).
\end{enumerate}
\end{solution}

\begin{problem}
某地为促进淡水鱼养殖业的发展，将价格控制在适当范围内，决定对淡水鱼养殖提供政府补贴. 设淡水鱼的市场价格为 \(x\) 元/千克，政府补贴为 \(t\) 元/千克. 根据市场调查，当 \(8 \leqslant x \leqslant 14\) 时，淡水鱼的市场日供应量 \(P\) 千克与市场日需求量 \(Q\) 千克近似地满足关系：\(P = 1000(x + t - 8)\)（\(x \geqslant 8\)，\(t \geqslant 0\)），\(Q = 500\sqrt{40 - (x - 8)^2}\)（\(8 \leqslant x \leqslant 14\)）.

当 \(P = Q\) 时，市场价格称为市场平衡价格.

\begin{enumerate}
\item 将市场平衡价格表示为政府补贴的函数，并求出函数的定义域.

\item 为使市场平衡价格不高于每千克 \(10\) 元，政府补贴至少为每千克多少元.
\end{enumerate}
\end{problem}

\begin{solution}
\begin{enumerate}
\item 令 \(u=x-8\)，则 \(0\leqslant u\leqslant6\). 由 \(P=Q\) 得
\[
1000(u+t)=500\sqrt{40-u^2}
\]
即 \(2(u+t)=\sqrt{40-u^2}\). 由于左边非负，可以平方，得到
\[
4(u+t)^2=40-u^2
\Longleftrightarrow 5u^2+8tu+4t^2-40=0
\]
解此关于 \(u\) 的方程，得
\[
u=-\frac45t\pm\frac25\sqrt{50-t^2}
\]
负号分支不满足 \(u\geqslant0\)，所以只能取正号. 又 \(t\geqslant0\)，且 \(u\geqslant0\) 要求
\[
\sqrt{50-t^2}\geqslant2t\Longleftrightarrow 0\leqslant t\leqslant\sqrt{10}
\]
在此范围内 \(u\leqslant2\sqrt2<6\)，所以 \(u\leqslant6\) 自动满足. 故市场平衡价格为
\(x(t)=8-\frac45t+\frac25\sqrt{50-t^2}\)，\(0\leqslant t\leqslant\sqrt{10}\).

\item 由上式，要求 \(x(t)\leqslant10\)，即
\[
8-\frac45t+\frac25\sqrt{50-t^2}\leqslant10
\Longleftrightarrow \sqrt{50-t^2}\leqslant2t+5
\]
两边均非负，可以平方，得
\[
50-t^2\leqslant(2t+5)^2
\Longleftrightarrow t^2+4t-5\geqslant0
\Longleftrightarrow (t-1)(t+5)\geqslant0
\]
结合定义域 \(t\geqslant0\)，得 \(t\geqslant1\). 因此政府补贴至少为每千克 \(1\) 元.
\end{enumerate}
\end{solution}

\begin{problem}
设 \(\{a_{n}\}\) 是由正数组成的等比数列，\(S_{n}\) 是其前 \(n\) 项和.

\begin{enumerate}
\item 证明：\(\frac{\lg S_n + \lg S_{n + 2}}{2} < \lg S_{n + 1}\).

\item 是否存在常数 \(c > 0\)，使得 \(\frac{\lg (S_n - c) + \lg (S_{n+2} - c)}{2} = \lg (S_{n+1} - c)\) 成立？并证明你的结论.
\end{enumerate}
\end{problem}

\begin{solution}
\begin{enumerate}
\item 设等比数列的首项为 \(a_1>0\)，公比为 \(q>0\). 当 \(q=1\) 时，\(S_n=na_1\)，于是
\[
S_nS_{n+2}-S_{n+1}^2
=na_1\cdot(n+2)a_1-(n+1)^2a_1^2
=-a_1^2<0
\]
当 \(q\ne1\) 时，\(S_n=a_1\dfrac{1-q^n}{1-q}\)，因此
\[
S_nS_{n+2}-S_{n+1}^2
=\frac{a_1^2}{(1-q)^2}\left[(1-q^n)(1-q^{n+2})-(1-q^{n+1})^2\right]
=-a_1^2q^n<0
\]
所以 \(S_nS_{n+2}<S_{n+1}^2\). 由于各项和均为正数，且 \(\lg\) 是增函数，
\[
\frac{\lg S_n+\lg S_{n+2}}2
=\lg\sqrt{S_nS_{n+2}}
<\lg S_{n+1}
\]

\item 若等式成立，则其两边的对数都有意义，且等价于
\[
(S_n-c)(S_{n+2}-c)=(S_{n+1}-c)^2
\]
利用 \(S_nS_{n+2}-S_{n+1}^2=-a_1^2q^n\) 以及 \(S_n+S_{n+2}-2S_{n+1}=a_{n+2}-a_{n+1}=a_1q^n(q-1)\)，左边减去右边为
\[
(S_n-c)(S_{n+2}-c)-(S_{n+1}-c)^2
=-a_1q^n\left[a_1-c(1-q)\right]
\]
若 \(q\geqslant1\)，上式严格为负，不可能等于零. 若 \(0<q<1\)，要使其为零必须 \(c=\dfrac{a_1}{1-q}\)，但此时
\[
S_n=\frac{a_1(1-q^n)}{1-q}<\frac{a_1}{1-q}=c
\]
所以 \(S_n-c<0\)，对数无意义. 两种情况均不可能，故不存在满足条件的常数 \(c>0\).
\end{enumerate}
\end{solution}

\begin{problem}
已知椭圆 \(\frac{x^2}{24} + \frac{y^2}{16} = 1\)，直线 \(l: \frac{x}{12} + \frac{y}{8} = 1\). \(P\) 是 \(l\) 上一点，射线 \(OP\) 交椭圆于点 \(R\)，又点 \(Q\) 在 \(OP\) 上且满足 \(|OQ| \cdot |OP| = |OR|^2\)，当点 \(P\) 在 \(l\) 上移动时，求点 \(Q\) 的轨迹方程，并说明轨迹是什么曲线.
\end{problem}

\begin{solution}
设 \(Q=(x,y)\)，以 \(OP\) 的方向为极轴方向，记单位方向向量为 \((\cos\theta,\sin\theta)\)，并设 \(OP=p\)、\(OQ=q\)、\(OR=r\)，其中 \(p\)，\(q\)，\(r>0\). 由于 \(P\in l\)，有
\[
p\left(\frac{\cos\theta}{12}+\frac{\sin\theta}{8}\right)=1
\]
由于 \(R\) 在椭圆上，有
\[
r^2\left(\frac{\cos^2\theta}{24}+\frac{\sin^2\theta}{16}\right)=1
\]
题设关系 \(OQ\cdot OP=OR^2\) 给出 \(q p=r^2\). 由前三式消去 \(p\)，\(r\)，得到
\[
q=\frac{\dfrac{\cos\theta}{12}+\dfrac{\sin\theta}{8}}
{\dfrac{\cos^2\theta}{24}+\dfrac{\sin^2\theta}{16}}
\]
又 \(x=q\cos\theta\)，\(y=q\sin\theta\)，所以 \(\cos\theta=\dfrac{x}{q}\)、\(\sin\theta=\dfrac{y}{q}\). 代入并约去正数 \(q\)，得
\[
\frac{x^2}{24}+\frac{y^2}{16}=\frac{x}{12}+\frac{y}{8}
\]
整理并配方：
\[
2x^2+3y^2=4x+6y
\Longleftrightarrow 2(x-1)^2+3(y-1)^2=5
\]
即
\[
\frac{(x-1)^2}{\frac52}+\frac{(y-1)^2}{\frac53}=1
\]
因为 \(q=|OQ|>0\)，所以 \(Q\ne O\)，即 \(x\)、\(y\) 不同时为零. 反过来，取上述椭圆上任意一点 \(Q\ne O\)，令 \(q=|OQ|>0\)，并令 \((\cos\theta,\sin\theta)=(x/q,y/q)\). 由椭圆方程可得 \(q\left(\frac{\cos^2\theta}{24}+\frac{\sin^2\theta}{16}\right)=\frac{\cos\theta}{12}+\frac{\sin\theta}{8}\)，所以右边为正. 令 \(p=\dfrac{1}{\frac{\cos\theta}{12}+\frac{\sin\theta}{8}}\)，\(r=\dfrac{1}{\sqrt{\frac{\cos^2\theta}{24}+\frac{\sin^2\theta}{16}}}\)，则 \(P=p(\cos\theta,\sin\theta)\) 在直线 \(l\) 上，\(R=r(\cos\theta,\sin\theta)\) 在原椭圆上，且由上式有 \(qp=r^2\). 因而每个这样的 \(Q\) 都能由题设构造得到，故轨迹正是上述椭圆去掉原点. 它的中心为 \((1,1)\)，长轴平行于 \(x\) 轴，长、短半轴分别为 \(\sqrt{\frac52}\)、\(\sqrt{\frac53}\).
\end{solution}
